\section{地址实现：一维比特带到三维局域访问}

\subsection{全局比特带与区间地址}

为实现“本体一维存储—三维局部访问”的结构，引入全局比特带
$$
B\in\{0,1\}^{\NN},
$$
并对层级节点集合 $V_\ell$ 定义区间地址映射
$$
\mathrm{Addr}_\ell:V_\ell\to\{I\subset\NN\},
$$
使得节点 $v$ 的局部本体读出由 $B[I(v)]$ 给出，再经 $\Fold_{m(v)}$ 与残差更新得到 $(x_v,r_v)$。该构造允许将空间邻域访问转化为比特带的局部切片访问。

\subsection{空间填充曲线与局域性增强（Hilbert 索引）}

空间填充曲线为高维网格到一维序列的双射（或近似双射）提供通用机制。Hilbert 曲线以其局域性性质（locality property）在索引与遍历任务中常被采用。\cite{wikipedia_hilbert_curve}

本文将 Hilbert 索引视为一种实现选择：其目标是使得在比特带上相近的地址区间更大概率对应于空间中相邻或近邻节点，从而减少跨越式通信并支持局部依赖图的高效调度。该选择不改变第 6--9 节的抽象定义，但影响第 11 节动态分辨率与第 15 节实验中对局部性成本的观测。\AuditTag 与第 3 节与第 9 节的口径一致，我们把“显示维数/显示字典”的确定视为 CAP 闭包输出；本文默认采用三维显示并在此以三维 Hilbert 索引作为一种可复核的地址实现。

\paragraph{Hilbert 显示的超立方体形状与自相似。}
当把长度为 $m$ 的地址（比特带上的有限前缀）映射到 $d$ 维显示空间时，一种标准做法是将地址位均匀分配到各坐标轴，使每轴携带
$$
n=\frac{m}{d}
$$
个有效位（在实现中对应对齐/填充与固定字典；当 $m$ 可被 $d$ 整除时最为自然）。因此显示网格的形状天然为 $d$ 维超立方体：
$$
\{0,1,\dots,2^n-1\}^d,
$$
即二维为 $2^n\times 2^n$ 的正方形、三维为 $2^n\times 2^n\times 2^n$ 的立方体。举例而言，$m=6$ 时二维对应 $n=3$（$8\times 8$），三维对应 $n=2$（$4^3$）；$m=12$ 时二维为 $64\times 64$，三维为 $16^3$。

此外，Hilbert 曲线按阶数递归构造：阶 $n$ 的曲线由 $2^d$ 个阶 $n-1$ 的子曲线经旋转/翻折后拼接而成，从而在几何上呈现显著的自相似结构。对本文而言，这一递归结构的意义在于：当 CAP 在不同显示维数 $d$ 之间权衡时，改变 $d$ 等价于在固定 $m$ 下改变每轴分辨率 $n=m/d$，从而在“更高维但更粗”与“更低维但更细”的多尺度族之间做可审计的选择。

\paragraph{从 $m$ 到 $m-d$ 的层级递归（自相似驱动的下一层级处理）。}
当 $n=m/d$ 足够大（例如三维显示下 $m=12$ 对应 $n=4$）时，Hilbert 的自相似允许将“在分辨率 $m$ 上的处理”递归地分解为“在下一层级 $m-d$ 上的同型处理”再加一个有限的拼接规则：整个 $d$ 维超立方体可分解为 $2^d$ 个边长减半的子超立方体
$$
\{0,1,\dots,2^n-1\}^d=\bigsqcup_{a\in\{0,1\}^d}\Bigl(a\cdot 2^{n-1}+\{0,1,\dots,2^{n-1}-1\}^d\Bigr),
$$
并且 Hilbert 顺序在每个子块内与阶 $n-1$ 的 Hilbert 顺序同构（仅差旋转/翻折与端点方向的有限状态）。因此，从一维地址的角度，高位自然决定子块索引 $a$（粗尺度），低位决定子块内坐标（细尺度）；从算法角度，任何依赖“局部块结构 + 有限边界条件”的计算（例如局部闭包、局部代价评估、或对 CAP 指标的分块统计）都可以在 $m$ 层通过对 $m-d$ 层的递归复用来实现。该递归视角与第 11 节的多尺度塔口径兼容：它为尺度投影 $\pi_{m\to m-d}$ 提供了一个具体、可复核且自相似一致的地址实现原型。

\subsection{Hilbert 递归尺度压缩作为截断族（遵循 Zeckendorf 码）}
\AuditTag 本小节把“尺度压缩 $m'\to m$”本身纳入定义 \ref{def:truncation_schema} 的可审计截断协议：空间输出仍强制落在稳定语言 $X_m$（Zeckendorf 编码），而尺度变化中被压缩/丢弃的部分以时间残差显式携带。

\paragraph{三类压缩规则（2D/3D/2D$\wedge$3D）。}
固定 $m<m'$。在实现层，我们用 Hilbert 索引为 $\Omega_{m'}$（地址长度为 $m'$ 的微观域）提供二维/三维的局域分块，并以此诱导一个从 $\Omega_{m'}$ 到 $\Omega_m$ 的粗粒候选集合（例如由“父块索引”给出）。随后对候选 $b\in\Omega_m$ 用第 5 节的折叠/规范化 $\Fold_m=\mathrm{Space}_m$ 投影到 $X_m$，得到候选空间态集合；最终的输出由截断族定义决定：
\begin{itemize}
  \item \textbf{2D 压缩：}以二维 Hilbert 的四叉分块递归为唯一局域性约束，输出空间态 $x\in X_m$ 与时间残差 $u$（记录块内偏移与必要的有限状态）。
  \item \textbf{3D 压缩：}以三维 Hilbert 的八叉分块递归为唯一局域性约束，输出空间态 $x\in X_m$ 与时间残差 $u$。
  \item \textbf{2D$\wedge$3D 同时压缩（两约束同时生效但不混合）：}二维与三维两套约束分别给出候选空间态集合；若两者存在公共可行候选，则在其上用 CAP 代价选择唯一代表；若发生冲突，则允许以 CAP 的“最小违约”策略选取空间代表，并把冲突证据/两套候选的差异作为时间残差的一部分显式写入 $u$。因此，两约束都参与输出，但不通过“混合显示”改变任一单独约束的定义。
\end{itemize}
在上述三族中，\emph{空间输出始终为 Zeckendorf 稳定词}（可审计规范形），而 Hilbert 的递归“朝向/翻折”有限状态、块内偏移序列以及尺度压缩引入的冲突证据均被归入时间残差（第 5 节的统一接口）。
